% Atelierin · Contrato de implantação e licença de uso
% Gerado em <<GERADO_EM>> a partir de /proposta/<<SLUG>> (escolhas da reunião).
% Compilação: XeLaTeX, duas vezes (sumário e total de páginas).
%   xelatex <<ARQUIVO>>

\documentclass[11pt,a4paper]{article}

% ----------------------------------------------------------------
% Pacotes e configuração de página
% ----------------------------------------------------------------
\usepackage[a4paper,left=2.35cm,right=2.35cm,top=2.35cm,bottom=2.5cm,headheight=15pt]{geometry}
\usepackage{fontspec}
\setmainfont{Inter}[
  UprightFont = *,
  BoldFont = * SemiBold,
  ItalicFont = * Italic,
  BoldItalicFont = * SemiBold Italic
]
\setsansfont{Inter}
\IfFontExistsTF{JetBrains Mono}{
  \setmonofont{JetBrains Mono}[Scale=0.86]
}{
  \setmonofont{DejaVu Sans Mono}[Scale=0.86]
}
\IfFontExistsTF{Bebas Neue}{
  \newfontfamily\displayfont{Bebas Neue}
}{
  \newfontfamily\displayfont{Inter}[BoldFont=Inter SemiBold]
}

\usepackage{microtype}
% No XeLaTeX o espaçamento de letras do microtype pode falhar; o do fontspec funciona sempre
\ifdefined\XeTeXversion\renewcommand{\textls}[2][100]{{\addfontfeatures{LetterSpace=\the\numexpr#1/10\relax}#2}}\fi
\usepackage{graphicx}
\usepackage{xcolor}
\usepackage[most]{tcolorbox}
\usepackage{enumitem}
\usepackage{booktabs}
\usepackage{tabularx}
\usepackage{array}
\usepackage{colortbl}
\usepackage{titlesec}
\usepackage{fancyhdr}
\usepackage{caption}
\usepackage[hidelinks]{hyperref}
\usepackage{float}
\usepackage{lastpage}
\usepackage{needspace}
\usepackage{setspace}

\hypersetup{
  pdftitle={Atelierin · Contrato de implantação e licença de uso · <<CLIENTE>>},
  pdfauthor={Samara Silvia Sabino},
  pdfsubject={Contrato do ambiente <<HOST>>}
}

\linespread{1.13}
\setlength{\parindent}{0pt}
\setlength{\parskip}{0.55em}
\setlength{\tabcolsep}{7pt}
\renewcommand{\arraystretch}{1.28}
\setlist[itemize]{leftmargin=1.35em,itemsep=0.25em,topsep=0.35em}
\setlist[enumerate]{leftmargin=1.55em,itemsep=0.3em,topsep=0.35em}
\setlength{\emergencystretch}{3em}
\hyphenpenalty=700
\exhyphenpenalty=700

% ----------------------------------------------------------------
% Identidade visual Atelierin — Crimson
% Capa dark + páginas internas light
% ----------------------------------------------------------------
\definecolor{atelierBg}{HTML}{080606}
\definecolor{atelierBgAlt}{HTML}{0F0D0D}
\definecolor{atelierCard}{HTML}{111010}
\definecolor{atelierSurface}{HTML}{1A1414}
\definecolor{atelierText}{HTML}{1A1414}
\definecolor{atelierTextSub}{HTML}{C4B8B0}
\definecolor{atelierMuted}{HTML}{806F69}
\definecolor{atelierDim}{HTML}{A99A93}
\definecolor{atelierAcc}{HTML}{C0211C}
\definecolor{atelierAccGlow}{HTML}{9E1A17}
\definecolor{atelierAccDim}{HTML}{4A1515}
\definecolor{atelierBorder}{HTML}{D9CEC7}
\definecolor{atelierBorderAcc}{HTML}{B65A54}
\definecolor{atelierPaper}{HTML}{F2EDE6}
\definecolor{atelierPaperCard}{HTML}{FDFAF7}
\definecolor{atelierGreen}{HTML}{477A55}
\definecolor{atelierAmber}{HTML}{A66B24}
\definecolor{atelierRed}{HTML}{A94444}
\arrayrulecolor{atelierBorder}

\color{atelierText}
\pagecolor{atelierPaper}

% Wordmark gráfico: o nome escrito continua sendo Atelierin.
\newcommand{\atelierwordmark}{%
  {\displayfont\fontsize{34}{36}\selectfont
  \textcolor{atelierPaperCard}{ATELIER}\textcolor{atelierAcc}{.IN}}}
\newcommand{\atelierwordmarklight}{%
  {\displayfont\fontsize{20}{22}\selectfont
  \textcolor{atelierText}{ATELIER}\textcolor{atelierAcc}{.IN}}}

% ----------------------------------------------------------------
% Cabeçalho, rodapé e títulos
% ----------------------------------------------------------------
\pagestyle{fancy}
\fancyhf{}
\fancyhead[L]{\small\color{atelierMuted}\textls[70]{ATELIER\textcolor{atelierAcc}{.IN}}}
\fancyhead[R]{\small\color{atelierMuted}Contrato \textperiodcentered\ <<HOST>>}
\renewcommand{\headrulewidth}{0pt}
\fancyfoot[L]{\small\color{atelierMuted}Rubricas: \rule{1.6cm}{0.4pt}\quad\rule{1.6cm}{0.4pt}}
\fancyfoot[R]{\small\color{atelierMuted}\thepage\ /\ \pageref{LastPage}}

\titleformat{\section}[display]
  {\Large\bfseries\color{atelierText}}
  {\small\color{atelierAcc}\textls[90]{CLÁUSULA \thesection}}{0.15em}{}
  [\vspace{-0.35em}{\color{atelierBorder}\rule{\linewidth}{0.7pt}}]
\titleformat{\subsection}
  {\large\bfseries\color{atelierAcc}}
  {\thesubsection}{0.65em}{}
\titleformat{\subsubsection}
  {\normalsize\bfseries\color{atelierText}}
  {\thesubsubsection}{0.55em}{}

\titlespacing*{\section}{0pt}{1.7em}{0.8em}
\titlespacing*{\subsection}{0pt}{1.25em}{0.4em}
\titlespacing*{\subsubsection}{0pt}{0.9em}{0.25em}

\captionsetup{font={small,color=atelierMuted},labelfont=bf,skip=6pt}

% ----------------------------------------------------------------
% Componentes visuais — páginas internas em light
% ----------------------------------------------------------------
\newtcolorbox{destaque}[1][]{
  enhanced,
  breakable,
  colback=atelierPaperCard,
  colframe=atelierBorderAcc,
  boxrule=0.65pt,
  borderline west={2pt}{0pt}{atelierAcc},
  arc=1mm,
  left=3.5mm,right=3.5mm,top=2.8mm,bottom=2.8mm,
  before skip=0.7em,after skip=0.8em,
  #1
}

\newtcolorbox{recomendacao}[1][]{
  enhanced,
  breakable,
  colback=atelierBgAlt,
  colframe=atelierBgAlt,
  coltext=atelierPaper,
  boxrule=0pt,
  borderline west={2.5pt}{0pt}{atelierAcc},
  arc=1mm,
  left=4mm,right=4mm,top=3.2mm,bottom=3.2mm,
  before skip=0.8em,after skip=0.9em,
  #1
}

\newtcolorbox{observacao}[1][]{
  enhanced,
  breakable,
  colback=atelierPaperCard,
  colframe=atelierBorder,
  boxrule=0.55pt,
  arc=1mm,
  left=3.2mm,right=3.2mm,top=2.5mm,bottom=2.5mm,
  before skip=0.6em,after skip=0.7em,
  #1
}

\newcommand{\rotulo}[1]{{\small\bfseries\color{atelierMuted}#1}\hspace{0.25em}}
\newcommand{\valor}[1]{{\bfseries\color{atelierAcc}#1}}
\newcommand{\tagrecomendada}{\tcbox[on line,colback=atelierAcc!9,colframe=atelierAcc!9,coltext=atelierAcc,boxrule=0pt,arc=1pt,left=5pt,right=5pt,top=2pt,bottom=2pt,fontupper=\footnotesize\bfseries]{RECOMENDADA}}
\newcommand{\tagestimativa}{\tcbox[on line,colback=atelierAmber!10,colframe=atelierAmber!10,coltext=atelierAmber,boxrule=0pt,arc=1pt,left=5pt,right=5pt,top=2pt,bottom=2pt,fontupper=\footnotesize\bfseries]{ESTIMATIVA}}
\newcommand{\opcao}[2]{\Needspace{4\baselineskip}\subsubsection*{#1\hfill #2}}
\newcommand{\campo}[1]{\textcolor{atelierMuted}{\textbf{#1}}}

\renewcommand{\contentsname}{Sumário}
\renewcommand{\figurename}{Figura}
\renewcommand{\tablename}{Tabela}

\newcommand{\linha}[1]{\rule{#1}{0.4pt}}

\begin{document}

% =================================================================
% CAPA
% =================================================================
\begin{titlepage}
  \thispagestyle{empty}
  \pagecolor{atelierBg}
  \color{atelierPaper}

  \vspace*{0.55cm}
  \noindent
  \begin{minipage}[t]{0.62\linewidth}
    \atelierwordmark
  \end{minipage}
  \begin{minipage}[t]{0.36\linewidth}
    \raggedleft
    {\fontspec{DejaVu Sans Mono}\fontsize{7.5}{10}\selectfont
    \color{atelierDim}\textls[90]{CONTRATO}\\[-1pt]
    \color{atelierAcc}\textls[90]{<<ANO>> \quad / \quad V. 1.0}}
  \end{minipage}

  \vspace{1.0cm}
  {\color{atelierAcc}\rule{0.19\linewidth}{2pt}}\par
  \vspace{1.25cm}

  {\fontspec{DejaVu Sans Mono}\fontsize{8}{11}\selectfont
  \color{atelierAcc}\textls[150]{ATELIERIN SCHOOL \quad / \quad AMBIENTE PERSONALIZADO}}\par

  \vspace{0.65cm}
  {\displayfont\fontsize{40}{52}\selectfont
  \color{atelierPaper}CONTRATO DE\\[0pt]
  \color{atelierPaper}IMPLANTAÇÃO\\[0pt]
  \color{atelierAcc}E LICENÇA DE USO}\par

  \vspace{0.5cm}
  {\Large\color{atelierTextSub}Ambiente digital com a sua marca,\\
  mantido pela Atelierin}\par

  \vspace{1.4cm}
  \noindent
  \begin{minipage}{0.94\linewidth}
    {\fontspec{DejaVu Sans Mono}\fontsize{7.5}{10}\selectfont
    \color{atelierDim}\textls[110]{AMBIENTE DIGITAL}}
    \par\vspace{0.2cm}
    {\Large\bfseries\color{atelierPaper}\texttt{<<HOST>>}}\par
  \end{minipage}

  \vfill
  \noindent\color{atelierAcc}\rule{\linewidth}{0.45pt}\par
  \vspace{0.55cm}

  \begin{tabularx}{\linewidth}{@{}>{\fontspec{DejaVu Sans Mono}\fontsize{7}{9}\selectfont\color{atelierDim}}p{0.22\linewidth} X@{}}
    CONTRATANTE & {\color{atelierPaper}\bfseries <<C_NOME>>}\\[0.45em]
    CONTRATADA & {\color{atelierPaper}\bfseries <<P_NOME>>}\\[0.45em]
    EMISSÃO & {\color{atelierPaper}<<DATA>>}\\
  \end{tabularx}

  \vspace{0.75cm}
  {\fontspec{DejaVu Sans Mono}\fontsize{7}{9}\selectfont
  \color{atelierDim}\textls[70]{ATELIERIN \quad / \quad CONTRATO DE IMPLANTAÇÃO E LICENÇA DE USO}}
\end{titlepage}

\pagecolor{atelierPaper}
\color{atelierText}
\pagestyle{fancy}

\tableofcontents
\newpage

% =================================================================
% PARTES
% =================================================================
{\large\bfseries Partes}\par
\vspace{0.4em}

\begin{tabularx}{\linewidth}{@{}>{\leavevmode\color{atelierMuted}}p{0.24\linewidth} X@{}}
  \toprule
  \multicolumn{2}{@{}l}{\color{atelierAcc}\bfseries CONTRATADA} \\
  Nome & <<P_NOME>> \\
  CPF ou CNPJ & <<P_DOC>> \\
  Endereço & <<P_END>> \\
  E-mail & <<P_EMAIL>> \\
  \midrule
  \multicolumn{2}{@{}l}{\color{atelierAcc}\bfseries CONTRATANTE} \\
  Nome & <<C_NOME>> \\
  CPF ou CNPJ & <<C_DOC>> \\
  Endereço & <<C_END>> \\
  E-mail & <<C_EMAIL>> \\
  \bottomrule
\end{tabularx}

As partes acima identificadas firmam este contrato, que reúne as explicações sobre o funcionamento da plataforma e as condições combinadas para a implantação e o uso do ambiente digital da Contratante, nos termos a seguir.

% =================================================================
\section{O que é o Atelierin}
% =================================================================

O Atelierin é uma plataforma digital criada, mantida e administrada pela Contratada, que oferece a profissionais que atendem pessoas (professores, consultores, terapeutas, entre outros) um ambiente próprio para organizar alunos ou clientes, cobranças, agenda, materiais e acompanhamento.

O sistema é o mesmo para todos os ambientes da plataforma. Cada ambiente recebe a identidade visual, o endereço, os módulos e o vocabulário do seu titular, e funciona separado dos demais, com os seus próprios dados.

\begin{destaque}
  \textbf{Natureza da contratação}\par
  A Contratante não compra um sistema. Ela contrata a montagem do seu ambiente e o direito de usá-lo enquanto este contrato estiver em vigor, com a parte técnica sob responsabilidade da Contratada.
\end{destaque}

% =================================================================
\section{Objeto}
% =================================================================

Este contrato tem por objeto:

\begin{enumerate}[label=\alph*)]
  \item a implantação do ambiente da Contratante no endereço \textbf{\texttt{<<HOST>>}};
  \item a licença de uso desse ambiente durante a vigência do contrato;
  \item a manutenção técnica, a hospedagem e o suporte do ambiente;
%<IF:basica>
  \item a criação de uma identidade básica para a Contratante;
%</IF:basica>
%<IFNOT:basica>
  \item a criação de um monograma para a Contratante, como cortesia;
%</IFNOT:basica>
  \item a página simples de apresentação no mesmo endereço<<OBJ_FIM>>
%<IF:lp>
  \item a criação de uma landing page completa para a Contratante.
%</IF:lp>
\end{enumerate}

% =================================================================
\section{O que entra na implantação}
% =================================================================

A implantação é o trabalho de montar o ambiente até ele estar no ar e testado, e compreende:

\begin{itemize}
  \item uma prévia com a marca da Contratante, apresentada antes do início;
  \item a aplicação da identidade visual: cores, fontes, logo ou monograma e tela de entrada;
  \item a escolha dos módulos que ficam ligados e do vocabulário da área da Contratante;
  \item o cadastro inicial dos alunos e do conteúdo que a Contratante já tiver, a partir dos materiais que ela fornecer;
  \item rodadas de ajuste até a aprovação da Contratante;
  \item uma conversa de entrega para mostrar como usar o ambiente.
\end{itemize}

%<IF:basica>
A identidade básica inclui logo simples, paleta de cores, tipografia, dois mockups com a marca aplicada e duas rodadas de ajuste. Não inclui manual de marca, papelaria ou peças para redes sociais. Se a Contratante quiser proteger o logo, o registro é feito no INPI, com a taxa paga diretamente por ela.
%</IF:basica>
%<IFNOT:basica>
O monograma é oferecido como \textbf{cortesia de primeiro cliente}, sem custo (valor de referência: <<MONOGRAMA>>). Inclui o nome da Contratante com uma tipografia escolhida, paleta de cores, um ícone com a inicial e uma rodada de ajuste. Por não ser um logo, não exige registro no INPI.
%</IFNOT:basica>

A página simples apresenta foto, nome, bio curta, serviços, contato e o botão de entrada no ambiente, a partir de um modelo pronto.
%<IF:lp>
A landing page completa tem texto, seções e design próprios, pensados para atrair novos alunos.
%</IF:lp>

% =================================================================
\section{Propriedade intelectual e licença de uso}
% =================================================================

O Atelierin, incluindo código-fonte, arquitetura, banco de dados, componentes, funcionalidades, identidade da plataforma e todas as melhorias feitas nele, pertence à Contratada. Isso vale também para o que for desenvolvido durante a implantação deste ambiente, que poderá ser usado em outros ambientes da plataforma.

A Contratante recebe uma licença de uso do seu ambiente, pessoal, intransferível e válida enquanto este contrato estiver vigente e os pagamentos estiverem em dia. O domínio \texttt{atelierin.com.br} e o subdomínio \texttt{<<HOST>>} continuam sob titularidade da Contratada.

%<IF:basica>
São da Contratante a sua marca, o seu conteúdo, os dados dos seus alunos e os valores que ela recebe. Os direitos sobre a identidade básica criada para ela são transferidos à Contratante após a quitação.
%</IF:basica>
%<IFNOT:basica>
São da Contratante a sua marca, o seu conteúdo, os dados dos seus alunos e os valores que ela recebe. O monograma criado para ela é de uso livre da Contratante.
%</IFNOT:basica>

Quem quiser um sistema exclusivo, com o código para si, deve contratar outro projeto, com proposta própria.

% =================================================================
\section{Valores e forma de pagamento da entrada}
% =================================================================

\begin{center}
\begin{tabularx}{0.92\linewidth}{@{}X r@{}}
  \toprule
  \textbf{Item} & \textbf{Valor} \\
  \midrule
<<ITENS>>
  \midrule
  \textbf{Total da entrada} & \textbf{<<TOTAL>>} \\
  \bottomrule
\end{tabularx}
\end{center}

<<PAGAMENTO>>

A implantação só começa após o pagamento da primeira parcela ou do valor à vista. O valor remunera o serviço de implantação e não representa a compra do sistema.

% =================================================================
\section{Licença e manutenção mensal}
% =================================================================

Além da entrada, a Contratante paga todo mês a licença e manutenção do ambiente, um valor fixo definido pela faixa de alunos ativos.

\begin{center}
\begin{tabularx}{0.92\linewidth}{@{}X r@{}}
  \toprule
  Faixa atual & <<FAIXA>> \\
  \textbf{Valor mensal} & \textbf{<<MENSAL>>} \\
  Vencimento & todo dia <<VENC>>, por Pix \\
  Primeiro pagamento & no mês seguinte à entrega do ambiente \\
  \bottomrule
\end{tabularx}
\end{center}

O valor inclui: licença de uso, hospedagem, banco de dados, domínio, cópias de segurança, atualizações de segurança, correções, monitoramento, suporte por mensagem e pequenos ajustes de texto e configuração. O valor é o mesmo havendo ou não pedidos de suporte no mês.

Se o número de alunos ativos mudar de faixa, o valor muda a partir do mês seguinte, com aviso prévio. As faixas da plataforma são: até 15 alunos, R\$ 190,00; de 16 a 30, R\$ 290,00; de 31 a 50, R\$ 390,00; acima disso, valor combinado entre as partes.

%<IF:lock>
O valor mensal fica \textbf{sem reajuste nos primeiros 12 (doze) meses}. Depois disso, pode ser reajustado uma vez por ano pelo IPCA acumulado do período, ou pelo índice que o substituir, com aviso de 30 dias.
%</IF:lock>
%<IFNOT:lock>
O valor mensal pode ser reajustado uma vez por ano, a contar da assinatura, pelo IPCA acumulado do período, ou pelo índice que o substituir, com aviso de 30 dias.
%</IFNOT:lock>

% =================================================================
\section{Infraestrutura e vídeos}
% =================================================================

Os serviços de infraestrutura usados pelo ambiente (banco de dados, armazenamento e hospedagem) são contratados e administrados pela Contratada e já estão incluídos no valor mensal. A Contratante não paga esses serviços à parte.

%<IF:bunny>
Os vídeos serão hospedados no Bunny Stream, contratado à parte. O custo estimado é de <<BUNNY>> por mês e pode variar conforme o armazenamento e o tráfego. Esse custo é repassado à Contratante pelo valor real da fatura, junto com a mensalidade.
%</IF:bunny>
%<IFNOT:bunny>
Os vídeos serão hospedados no YouTube ou no Google Drive da Contratante e incorporados ao ambiente, sem custo adicional, sujeitos às regras e configurações de privacidade desses serviços. A troca para hospedagem própria (Bunny Stream) pode ser combinada depois, com custo à parte.
%</IFNOT:bunny>

% =================================================================
\section{Pagamento dos alunos}
% =================================================================

Os alunos pagam diretamente à Contratante, pela chave Pix ou conta que ela indicar. O dinheiro não passa pela Contratada nem pela plataforma.

A Contratante confere o recebimento na própria conta e marca o aluno como pago no painel financeiro do ambiente, onde acompanha quem está em dia, pendente ou atrasado. O envio de um comprovante pelo aluno não substitui essa conferência.

Uma integração bancária automática (por exemplo, cobrança por Pix gerada pelo sistema) pode ser contratada depois, com orçamento próprio, e depende das condições e tarifas do banco ou provedor escolhido.

% =================================================================
\section{Etapas e prazos}
% =================================================================

\begin{enumerate}
  \item \textbf{Conversa inicial} e envio dos materiais pela Contratante (logo ou referências, lista de alunos, conteúdos).
  \item \textbf{Aprovação da prévia} com a marca aplicada.
  \item \textbf{Entrega} do ambiente publicado, com uma conversa para mostrar como usar.
\end{enumerate}

\begin{center}
\begin{tabularx}{0.92\linewidth}{@{}X r@{}}
  \toprule
  Ambiente Personalizado & de 1 a 4 semanas \\
  <<IDENT_NOME>> & <<IDENT_PRAZO>> \\
%<IF:lp>
  Landing page & 3 dias úteis, podendo chegar a 5 \\
%</IF:lp>
  \bottomrule
\end{tabularx}
\end{center}

Os prazos contam a partir do pagamento da entrada e do recebimento dos materiais. Atrasos no envio de materiais ou nas aprovações da Contratante estendem os prazos pelo mesmo período.

% =================================================================
\section{Obrigações da Contratada}
% =================================================================

\begin{itemize}
  \item implantar, configurar e manter o ambiente conforme este contrato;
  \item administrar a infraestrutura técnica e manter cópias de segurança;
  \item corrigir falhas e prestar suporte por mensagem em horário comercial;
  \item avisar com antecedência sobre mudanças relevantes no ambiente, na faixa ou nos valores;
  \item tratar os dados do ambiente apenas para fazê-lo funcionar, com medidas razoáveis de segurança.
\end{itemize}

% =================================================================
\section{Obrigações da Contratante}
% =================================================================

\begin{itemize}
  \item fornecer os materiais e as informações necessários à implantação, nos prazos combinados;
  \item pagar a entrada e a mensalidade nas datas combinadas;
  \item conferir os pagamentos dos alunos e mantê-los atualizados no painel;
  \item cuidar da relação comercial com os seus alunos, incluindo valores, regras e comunicação;
  \item informar à Contratada as mudanças de faixa de alunos e os problemas encontrados no ambiente;
  \item usar o ambiente de acordo com a sua finalidade e com a lei.
\end{itemize}

% =================================================================
\section{Suporte, escopo e novas funcionalidades}
% =================================================================

O suporte mensal cobre o funcionamento do ambiente, correções e pequenos ajustes de texto e configuração. Não acumula de um mês para o outro.

Não estão incluídos e dependem de orçamento aprovado antes:

\begin{itemize}
  \item funcionalidades novas ou feitas só para a Contratante;
  \item integrações com bancos e outros sistemas;
  \item produção de conteúdo (aulas, textos, fotos);
  \item mudanças estruturais no fluxo de trabalho.
\end{itemize}

Nada é considerado contratado sem acordo prévio por escrito, inclusive por mensagem.

% =================================================================
\section{Disponibilidade e serviços de terceiros}
% =================================================================

O ambiente depende de serviços de terceiros (hospedagem, banco de dados, provedores de e-mail e de vídeo). A Contratada administra esses serviços e age para resolver problemas o quanto antes, mas não responde por interrupções causadas por esses fornecedores, falhas de rede ou eventos fora do seu controle.

% =================================================================
\section{Proteção de dados e confidencialidade}
% =================================================================

As partes seguem a Lei nº 13.709/2018, Lei Geral de Proteção de Dados Pessoais (LGPD). A Contratante é a controladora dos dados dos seus alunos e decide como eles são usados. A Contratada atua como operadora e trata esses dados somente para manter o ambiente funcionando, sem vendê-los ou compartilhá-los para outros fins.

A Contratante deve coletar apenas os dados necessários e informar os seus alunos sobre o uso do ambiente. As duas partes mantêm em sigilo as informações comerciais, técnicas e pessoais a que tiverem acesso por causa deste contrato.

% =================================================================
\section{Atraso no pagamento}
% =================================================================

Pagamentos em atraso têm multa de 2\% e juros de 1\% ao mês, proporcionais aos dias de atraso. Com mais de 15 dias de atraso, e após aviso, o acesso ao ambiente pode ser suspenso até a regularização, sem perda dos dados.

% =================================================================
\section{Vigência, cancelamento e encerramento}
% =================================================================

Este contrato vale por prazo indeterminado, a partir da assinatura. Qualquer parte pode encerrá-lo com aviso por escrito de 30 (trinta) dias, inclusive por e-mail ou mensagem.

\begin{itemize}
  \item valores vencidos e parcelas da entrada já devidas continuam devidos;
  \item a mensalidade do período de aviso é paga normalmente;
  \item antes do encerramento, a Contratante pode exportar os dados dos alunos e o conteúdo do ambiente;
  \item depois do encerramento, a licença de uso termina e os dados são apagados em até 60 dias, salvo obrigação legal de guarda;
  \item o sistema, o domínio e o subdomínio continuam da Contratada.
\end{itemize}

% =================================================================
\section{Disposições gerais}
% =================================================================

Alterações de escopo, valores ou condições só valem se combinadas por escrito entre as partes. As partes aceitam a assinatura eletrônica deste contrato. Fica eleito o foro da comarca de \textbf{Recife, Pernambuco}, para resolver qualquer questão sobre este contrato.

% =================================================================
\clearpage
{\large\bfseries Assinaturas}\par
% =================================================================

\vspace{0.6em}
E, por estarem de acordo, as partes assinam este contrato.

\vspace{0.8em}
Recife, <<DATA>>.

\vspace{2cm}

\begin{tabularx}{\linewidth}{@{}X X@{}}
  \rule{0.92\linewidth}{0.5pt} & \rule{0.92\linewidth}{0.5pt} \\
  \textbf{Contratada} & \textbf{Contratante} \\
  <<P_NOME>> & <<C_NOME>> \\
  <<P_DOC>> & <<C_DOC>> \\
\end{tabularx}

\vfill
\begin{center}
  {\color{atelierAcc}\rule{0.18\linewidth}{1.5pt}}\par
  \vspace{0.45em}
  {\atelierwordmarklight}\par
  {\small\color{atelierMuted}Um ambiente próprio para cada projeto.}
\end{center}

\end{document}
